\scanpage{20}
\mspage{18}

qu'il faut comparer avec l'élément déduit \unsure{du|des} \unsure{chemin|chemins} \unsure{tronqué|tronqués}

\diagram{\begin{tikzcd}[ampersand replacement=\&,column sep=small,row sep=tiny]
a\arrow[r]\&0\arrow[d,equal]\&b\arrow[l]\&\cdots\\
\&a_0
\end{tikzcd}}

Il faudrait montrer que $0$ est aussi un objet neutre de $M\Sigma^{-1}$, i.e. que

\[
 \deleted{\unsure{\mathrm{Hom}(\,}}\ \unsure{\Pi|\pi}_0(a,0)=\pi_0(0,a)=\{0\}\qquad\text{pour }\deleted{\text{tout}}\ a\in\mathrm{Ob}\,M,
\]

et interpréter \deleted{\unsure{les}} l'élément nul dans $\pi_0(a,b)$ comme l'application nulle dans $M\Sigma^{-1}$. \deleted{\unsure{Il}} En général, il faut des conditions qui assurent (sans que $M$ soit forcément \unsure{ponctuelle|ponctuable}) que

\[
 M\mathbin{\overset{\gamma}{\unsure{-|\longrightarrow}}}M\Sigma^{-1}
\]

commute aux sommes finies, \deleted{\unsure{ou}} aux produits finis (avec $\gamma(\unsure{\phi|\varnothing}_M)=\phi_{M\Sigma^{-1}}$, resp. $\gamma(1_M)=e_{M\Sigma^{-1}}$).

% slanted
\margin{à faire}

Supposons enfin que $M$ soit \uline{additive}, et \deleted{\unsure{soit $k$}} \unsure{même} $k$-linéaire ($k$ un anneau fixé), i.e. on a de plus un hom d'anneaux

\[
 k\longrightarrow\mathrm{End}(\mathrm{id}_M).
\]

On doit supposer quelque compatibilité de $\Sigma$ à la structure additive, p.ex. que la somme directe de deux flèches de $\Sigma$ est dans $\Sigma$, et \deleted{\unsure{de plus}} sans doute qque chose de plus fort, du genre ``lemme des \unsure{cinq}'' pour $\Sigma$. (Mais on ne suppose pas forcément $M$ abélienne donc on aurait du mal à \unsure{formuler|prétendre} un lemme des cinq ?)
% archive sheet number: 19
